\begin{abstract}
Self-improving agents increasingly convert terminal feedback into persistent memory, yet common evaluations conflate whether feedback changes stored state, whether that state is retrieved, whether the policy uses the branch-specific difference, and whether behavior changes. We freeze each source trajectory, switch only the released success/failure writer branch, and measure write $W$, native source-item exposure $E$, first-action uptake $U$, terminal outcome $O$, plus a forced-capacity control $F$. Writing diverges for all 20 Shopping and 4 Reddit pairs; forced supply has downstream leverage ($|\Delta|=0.15625$, $p=0.00074$). Under native Shopping reuse, however, source memory is retrieved in 125/172 opportunities while first-action uptake is unsupported (TV $=0.06944$, $p=0.5801$, 0/36 modal changes) and terminal transport is sparse ($|\Delta|=0.02083$, $p=0.4289$, 34/36 zero). The prerequisite order $O\Rightarrow U\Rightarrow E\Rightarrow W$ with independent $F$ induces a 10-state diagnostic class; exhaustive enumeration of all $2^5$ observation subsets shows that only the full $W/E/U/O/F$ basis separates every state. We then test repair actionability prospectively. A matched explicit memory-use reminder has no aggregate effect on 13 outcome-blind states (312/312 calls; $U_{A0}=U_{A1}=U_{A2}=0.0962$). On a separate fresh 12-state pilot, an intentionally non-novel state-conditioned binding adapter raises mean uptake from $0.125$ natively and $0.083$ under a memory-only adapter to $0.208$; however, the gain is heterogeneous (mean $U_{A2}-U_{A1}=0.125$, positive in only 3/12 states), so its preregistered consistency gate fails. Thus stage resolution can identify a repair family worth probing and produce a conditional transport opening, but it does not by itself supply a universal repair prescription. Persistent-state divergence is not behavioral authority, availability is not use, and repair authority requires separate intervention evidence.
\end{abstract}
